\domainevide{Les cotes de ce thème supposent le site étale d'un schéma et sa cohomologie, que mathlib ne contient pas.}{31, 33, 60}
\noindent Pistes de l'onglet Findings pour ces cotes, hors de portée de Lean pour l'instant :
\begin{itemize}\small\setlength{\itemsep}{0pt}
\item Cote~31 : une fonction $\varphi$ sur les points d'un schéma noethérien qui vérifie deux conditions locales majore la dimension cohomologique étale $\ell$-primaire. Mathlib n'a pas la cohomologie étale.\piste{31-phi-specialisation-criterion}
\item Cote~31 : conjecture : en la caractéristique résiduelle $\ell$, la borne de la dimension cohomologique d'un schéma affine est gouvernée par le degré d'imperfection $\nu(k)$, défini par $\ell^{\nu(k)} = [k : k^\ell]$. Mathlib n'a pas la dimension cohomologique étale.\piste{31-imperfection-bound-conjectures}
\item Cote~31 : pour un anneau local noethérien intègre $A$ de corps des fractions $K$ et de corps résiduel $k$, $\mathrm{cd}_r(K) \geq \mathrm{cd}_r(k) + \dim A$ pour tout premier $r$, à deux exceptions près. Mathlib n'a pas la dimension cohomologique galoisienne.\piste{31-local-ring-cd-lower-bound}
\item Cote~60 : pour un faisceau $\ell$-adique au point générique d'une courbe sur $\mathbf F_p$, la moyenne sur l'inertie donne un faisceau virtuel $E_\eta^\natural$ dont la fonction $L$ est additive et locale ; le dossier calcule son dual de Verdier. Mathlib n'a pas les faisceaux $\ell$-adiques, leurs fonctions $L$ ni la dualité de Verdier.\piste{60-inertia-average-duality}
\end{itemize}
